\documentclass{article}
\usepackage{geometry, amsmath, amsfonts}
\geometry{margin=1in, top=0.8in}

\begin{document}

\noindent\textbf{Schrodinger's Wave Equation For a Hydrogen Atom - By Leo}\\[0.2em]
\noindent\rule{\textwidth}{0.4pt}
\textit{This was written while I was in Year 12 as a project to help me learn all 4 of my subjects further and with greater rigour - using AI to teach me (not to fully write) the physics and how to solve these equations so that I understand these equations further than just regurgitating definitions. Im writing this from the perspective of teaching a former self these concepts so that I can consolidate the why and how. Although I wouldn't be able to solve this on my own I wanted to understand the why and not just use the maths without knowing what it did}
\section{Parametric 3D space}
Any position in 3D space can be written parametrically in the form of a spherical co-ordinate:

\begin{align*}
x = r\sin\theta\cos\phi, \quad
y = r\sin\theta\sin\phi, \quad
z = r\cos\theta 
\end{align*}
\begin{align*}
\text{with }   r& \geq0,\quad \theta \in [0, \pi],\quad \phi \in [0,2\pi),\quad \text{ and volume element } dV = r^2\sin\theta dr d\theta d\phi
\end{align*}
Here r is the distance from the origin to a point in 3D space, $\theta$ is the polar angle, measured from $+z$, and $\phi$ is the azimuthal angle in the $xy$ plane, measured from $+x$.\\\\
The volume element is the factor you need to multiply to get the new volume when changing the volume, for cartesian co-ordinates this is trivial, $dV = dxdydz$, as you are increasing the size of a cuboid, but with spherical co-ordinates this becomes a bit tricky as you aren't increasing the size of a box.\\\\
As a sanity check you could integrate across it:
\begin{align*}
V &= \int_0^R \int_0^\pi \int_0^{2\pi} r^2\sin\theta \,d\phi \,d\theta \,dr\\
&= 2\pi \int_0^R \int_0^\pi r^2\sin\theta \,d\theta \,dr\\
&= 4\pi \int_0^R r^2 \,dr\\
&= \frac{4\pi}{3} R^3
\end{align*}
which is exactly what you would expect from a sphere.\\\\

\section{The Laplacian}
The Laplacian is a differential operator, the Laplacian takes function $f(x,y,z)$ and gives a measure of how much the function curves at each point.\\
For cartesian co-ordinates the Laplacian is the sum of the second partial derivatives with respect to each co-ordinates, written like this:
\begin{align*}
\nabla^2f = \frac{\partial^2f}{\partial x^2} + \frac{\partial^2f}{\partial y^2} + \frac{\partial^2f}{\partial z^2}        
\end{align*}
So for example if you had $f(x,y,z)=x^2 + y^2 + z^2$
\begin{align*}
    \frac{\partial^2f}{\partial x^2} =2, \quad \frac{\partial^2f}{\partial y^2} = 2, \quad \frac{\partial^2f}{\partial z^2} = 2
    &\implies\nabla^2f = 6
\end{align*}
\newpage
\noindent Now for spherical co-ordinates the Laplacian is:
\begin{align*}
      \nabla^2f
  = \frac{1}{r^2}\frac{\partial}{\partial r}\!\left(r^2 \frac{\partial}{\partial r}\right)
  + \frac{1}{r^2\sin\theta}\frac{\partial}{\partial\theta}\!\left(\sin\theta \frac{\partial}{\partial\theta}\right)
  + \frac{1}{r^2\sin^2\theta}\frac{\partial^2}{\partial\phi^2}.
\end{align*}
Which you can derive from doing the Laplacian of $f(x(r,\theta,\phi),\,y(r,\theta,\phi),\,y(r,\theta,\phi))$ using the chain rule.

\section{The Schrödinger Equation}
\begin{itemize}
  \item A hydrogen atom is a proton with one electron bound to it by electrostatic attraction.
  \item Classical physics says that the electron should spiral into the proton (accelerating charge radiates energy). It does not, so the classical physics is wrong at this scale.
  \item Quantum mechanics replaces "the electron is at position $r$" with "the electron's state is described by a wavefunction $\psi(r)$."
  \item The wavefunction itself isn't directly observable. What is observable is $|\psi(r)|^2$, which gives the probability density of finding the electron at r.
  \item An orbital is a particular allowed wavefunction, labelled by quantum numbers.
\end{itemize}
These are adopted postulates from quantum mechanics for our equation, no further assumptions are required.\\\\
The Schrödinger Equation is stated as:
\begin{align*}
    - \frac{\hbar^2}{2m_e}
    \nabla^2\psi(r)+
    V(r)\psi(r) =
    E\psi(r)
\end{align*}

Unpacking this term by term:
\begin{itemize}
\item $- \frac{\hbar^2}{2m_e}$ is the kinetic energy term, the Laplacian measures the curvature of the wave function - a wave function with more curvature costs more kinetic energy.
\item $V(r)\psi$ is the potential energy term. For hydrogen, V is the Coulomb attraction between the electron and proton.
\item $E\psi$ is the total energy, where E is the eigenvalue and $\psi$ is the corresponding wavefunction (eigenfunction).
\end{itemize}
The constants $\hbar,m_e,e,4\pi\epsilon_0$ clutter our equations, and therefore we adopt atomic units, where\\ $\hbar = m_e = e = 4\pi\epsilon_0 = 1$ in these units, distances are measured in Bohr radii and energies in hartrees. After these changes our equation simplifies dramatically.\\
The new equation (given in atomic units) is given as:
\begin{align*}
    -\frac{1}{2} \nabla^2\psi + V\psi = E\psi
 \end{align*}
\newpage
 \section{The Coulomb Potential}
 The coulomb potential for hydrogen (in atomic units) is:
 \begin{align*}
     V(r) = -\frac{1}{|r|} = -\frac{1}{r}
 \end{align*}
 Here the minus sign means the attraction to the nucleus. The potential depends only on the distance r from the nucleus, not the direction, we can exploit this property later to seperate the equation.

 \section{The Equation To Be Solved}
 We now finally have the equation:
 \begin{align*}
     -\frac{1}{2} \nabla^2\psi - \frac{1}{r}\psi(r) = E\psi
 \end{align*}
 Which has the following boundary conditions:
 \begin{itemize}
     \item Normalisability: $\int |\psi|^2dV = 1$, so $|\psi|^2$ is a probability distribution
     \item Vanishing at infinity: $\psi \rightarrow 0 \text{ as }r \rightarrow \infty$
     \item Single-valued: $\psi$ has at most one value at each point
 \end{itemize}
 \section{Manipulating The Equations To Solvable Forms}
  \subsection{Substituting}
 Before we noted that the potential $V(r) = -\frac{1}{r}$ depends only on the distance from the nucleus, not on direction. The atom looks the same from every angle. We will now try a wavefunction of the form $\psi(r,\theta,\phi) = R(r)\Theta(\theta)\Phi(\phi)$\\\\
 substituting $\psi(r,\theta,\phi) = R(r)\Theta(\theta)\Phi(\phi)$ into our previous Schrödinger equation we recieve:
\[
      -\frac{1}{2}\left[
    \frac{1}{r^2}\frac{\partial}{\partial r}\!\left(r^2\frac{\partial\psi}{\partial r}\right)
    + \frac{1}{r^2\sin\theta}\frac{\partial}{\partial\theta}\!\left(\sin\theta\frac{\partial\psi}{\partial\theta}\right)
    + \frac{1}{r^2\sin^2\theta}\frac{\partial^2\psi}{\partial\phi^2}
  \right]
  - \frac{1}{r}\psi = E\psi\\
\]
subsituting $\psi = R\Theta\Phi$
\begin{align*}
    \frac{\partial}{\partial r} (r^2 \frac{\partial(R\Theta\Phi)}{\partial r}) &= \Theta\Phi\frac{d}{dr}(r^2\frac{dR}{dr})\\\\
    \frac{\partial}{\partial \theta} (\sin\theta \frac{\partial(R\Theta\Phi)}{\partial r}) &= R\Phi\frac{d}{d\theta}(\sin\theta\frac{d\Theta}{d\theta})\\\\
    \frac{\partial^2(R\Theta\Phi)}{\partial\phi^2} &= R\Theta\frac{d^2\Phi}{d\phi^2} 
\end{align*}
Ordinary d's replace partials as each function now depends only on one variable. Substituting again:
\[
  -\frac{1}{2}\left[
    \Theta\Phi\,\frac{1}{r^2}\frac{d}{dr}\!\left(r^2\frac{dR}{dr}\right)
    + \frac{R\Phi}{r^2\sin\theta}\frac{d}{d\theta}\!\left(\sin\theta\frac{d\Theta}{d\theta}\right)
    + \frac{R\Theta}{r^2\sin^2\theta}\frac{d^2\Phi}{d\phi^2}
  \right]
  - \frac{1}{r}R\Theta\Phi
  = E\,R\Theta\Phi
\]
\newpage
Divide through by $R\Theta\Phi$
\[
  -\frac{1}{2}\left[
    \frac{1}{R}\frac{1}{r^2}\frac{d}{dr}\!\left(r^2\frac{dR}{dr}\right)
    + \frac{1}{\Theta}\frac{1}{r^2\sin\theta}\frac{d}{d\theta}\!\left(\sin\theta\frac{d\Theta}{d\theta}\right)
    + \frac{1}{\Phi}\frac{1}{r^2\sin^2\theta}\frac{d^2\Phi}{d\phi^2}
  \right]
  - \frac{1}{r} = E
\]
Now multiply through by $-2r^2$ and rearrange
\[
  \frac{1}{R}\frac{d}{dr}\!\left(r^2\frac{dR}{dr}\right)
  + 2r^2\left(E + \frac{1}{r}\right)
  + \frac{1}{\Theta\sin\theta}\frac{d}{d\theta}\!\left(\sin\theta\frac{d\Theta}{d\theta}\right)
  + \frac{1}{\Phi\sin^2\theta}\frac{d^2\Phi}{d\phi^2}
  = 0
\]
\subsection{Seperating r From The Angles}
The equation must hold for all points in space. Consider varying r while holding $\theta$ and $\phi$ constant. the first group of terms only vary as r varies, while the further terms are unaffected by r. So for their sum to equal zero regardless of $r$, each group must seperately be constant - and those constants must be equal and opposite

\[
  \underbrace{\frac{1}{R}\frac{d}{dr}\!\left(r^2\frac{dR}{dr}\right)+ 2r^2\left(E + \frac{1}{r}\right)}_{\text{depends only on r}}
  + 
  \underbrace{\frac{1}{\Theta\sin\theta}\frac{d}{d\theta}\!\left(\sin\theta\frac{d\Theta}{d\theta}\right)+ \frac{1}{\Phi\sin^2\theta}\frac{d^2\Phi}{d\phi^2}}_{\text{depends only on $\theta\phi$}}
  = 0
\]
Introducing the seperation constant $\lambda$:
\begin{align*}
\frac{1}{R}\frac{d}{dr}\!\left(r^2\frac{dR}{dr}\right)+ 2r^2\left(E + \frac{1}{r}\right) &= \lambda\\\\
\frac{1}{\Theta\sin\theta}\frac{d}{d\theta}\!\left(\sin\theta\frac{d\Theta}{d\theta}\right)+ \frac{1}{\Phi\sin^2\theta}\frac{d^2\Phi}{d\phi^2} &= -\lambda
\end{align*}
\subsection{Second Seperation of $\theta$ from $\phi$}
Multiply the second equation by $\sin^2\theta$:
\[
\frac{\sin\theta}{\Theta}\frac{d}{d\theta}(\sin\theta\frac{d\Theta}{d\theta}) + \lambda\sin^2\theta + \frac{1}{\Phi}\frac{d^2\Phi}{d\phi^2} =0 
\]
Now again we see that the first terms depend only on $\theta$ while the further terms only depend on $\phi$. So we can introduce a second seperation constant. We write this as $-m^2$, we write this constant as $-m^2$ as the equation for $\Phi$ will have oscillatory solutions in the form $e^{\pm im\phi}$.\\
Then:
\[
\frac{1}{\Phi}\frac{d^2\Phi}{d\phi^2} =-m^2
\]
\[
\frac{\sin\theta}{\Theta}\frac{d}{d\theta}(\sin\theta\frac{d\Theta}{d\theta}) + \lambda\sin^2\theta = m^2
\]
Re-arranging both into solvable forms:
\[
\frac{d^2\Phi}{d^2\phi} = -m^2
\]
\[
\frac{1}{\sin\theta}\frac{d}{d\theta}\!\left(\sin\theta\frac{d\Theta}{d\theta}\right)
  + \left(\lambda - \frac{m^2}{\sin^2\theta}\right)\Theta = 0
\]
\newpage
\subsection{Final Solvable Forms}
\[ (R) \quad\frac{1}{R}\frac{d}{dr}\!\left(r^2\frac{dR}{dr}\right)+ 2r^2\left(E + \frac{1}{r}\right) = \lambda \]
\[ (\Theta) \quad \frac{1}{\sin\theta}\frac{d}{d\theta}\!\left(\sin\theta\frac{d\Theta}{d\theta}\right)
  + \left(\lambda - \frac{m^2}{\sin^2\theta}\right)\Theta = 0\]
\[(\Phi) \quad \frac{d^2\Phi}{d^2\phi} = -m^2\]

\section{Solving These New Equations}
\subsection{First ODE}
Out of these equations the simplest to solve will be the azimuthal equation for $\Phi(\phi)$\\
Assume $\Phi=e^{k\phi}$ and differentiate twice:
\[
\frac{d\Phi}{d\phi} =ke^{k\phi}, \quad \frac{d^2\Phi}{d\phi^2} =k^2e^{k\phi}\\
\]
\[
k^2e^{k\phi} = -m^2e^{k\phi}
\]
\[
k^2 = -m^2 \implies k=\pm im
\]
\[
 \Phi = Ae^{im\phi} + Be^{-im\phi}
\]
Now we must impose boundary conditions. The wave function must be single valued meaning that a point in space cannot have two different probabilities. Since $\phi$ and $2\pi + \phi$ describe the same physical point we require:
\[
\Phi(\phi + 2\pi) = \Phi (\phi)
\]
Using $\Phi(\phi)=e^{im\phi}$:
\[
e^{im(\phi +2\pi)} = e^{im\phi} \implies e^{2\pi im} = 1
\]
\[m =0,\pm1,\pm2,...\]
\[
    \boxed{m\in \mathbb{Z}}
\]
\subsection{Normalising This Result}
Now we need to fix the constant A by requiring
\[
\int^{2\pi}_0 |\Phi|^2d\phi = 1
\]
For the complex form $\Phi = Ae^{im\phi}$, this is:
\[
\int^{2\pi}_0 |A|^2d\phi = 2\pi|A|^2 = 1 \implies |A| = \frac{1}{\sqrt{2\pi}}
\]
Giving:
\[
\boxed{\Phi_m(\phi) = \frac{1}{\sqrt{2\pi}}e^{im\phi}, \quad m\in \mathbb{Z}}
\]
So we have recieved our first quantum number, m, which controls how the wavefunctions varies as you go around the z axis.
\newpage
\subsection{The Second ODE}
The first azimuthal equation gave us our first quantum number m. Now we need to solve the polar equation $\Theta(\theta)$ which turns out being much harder.\\
We need to solve:
\[
\quad \frac{1}{\sin\theta}\frac{d}{d\theta}\!\left(\sin\theta\frac{d\Theta}{d\theta}\right)
  + \left(\lambda - \frac{m^2}{\sin^2\theta}\right)\Theta = 0
  \]
With $\lambda$ still an undertemined seperation constant meaning, unlike the previous ODE, this has non constant coefficients, and its solutions are not elementary functions. The equation involves both $\sin\theta$ and $\cos\theta$, and derivatives with respect to $\theta$. If we substitute $u = \cos\theta$, trigonometric functions become polynomials which are easier to work with.

\[
  \frac{d}{d\theta}
  = \frac{du}{d\theta}\frac{d}{du}
  = -\sin\theta\frac{d}{du}
\]
So
\[
\sin\theta\frac{d}{d\theta} = \sin\theta(-\sin\theta\frac{d}{du}) = -\sin^2\theta\frac{d}{d\theta} = -(1-u^2)\frac{d}{du}
\]
And applying twice like the equation does
\[
\sin\theta\frac{d}{d\theta}(\sin\theta\frac{d\Theta}{d\theta}) = (1-u^2)\frac{d}{du}[(1-u^2)\frac{dP}{du}]
\]
Where $P(u) = \theta(\arccos u)$
Substitute into the equation, divide through by $\sin^2\theta$ and you get:
\[
  \frac{d}{du}\!\left[(1-u^2)\frac{dP}{du}\right]
  + \left(\lambda - \frac{m^2}{1-u^2}\right)P = 0
\]
This is the associated Legrendre equation which is a standard second-order ODE which solutions are known and we will use.
\end{document}


